\documentclass{beamer}

\usepackage[utf8]{inputenc}
\usepackage[spanish]{babel}
\usepackage[T1]{fontenc}
\usepackage{graphicx}
\usepackage{tikz}

\renewcommand\shorthandsspanish{}
\noextrasspanish

\title{Python}
\author{
Guillem Borrell i Nogueras\\
\texttt{guillemborrell@gmail.com}
}

\begin{document}

\begin{frame}
\begin{center}
 \includegraphics[width=9cm]{files/python-logo-generic.pdf}\\
 % python-logo-generic.pdf: 389x115 pixel, 72dpi, 13.72x4.06 cm, bb=0 0 389 115
\begin{large}
\textbf{Python en Supercomputación}
\end{large}\\
Charla Técnica\\

Guillem Borrell i Nogueras\\

ETSIA, Octubre 2007
\end{center}

\end{frame}

\begin{frame}
 \frametitle{Charla técnica}
\begin{itemize}
 \item Python Vs. Matlab.
 \item Más sobre los wrappers.
 \item Arrays en C, Fortran y Python.
 \item ctypes.
 \item F2Py.
 \item Python en paralelo.
 \item GIL.
\end{itemize}

\end{frame}

\begin{frame}
 \frametitle{Python vs. Matlab}
\begin{center}
 ¿Ayuda Matlab a manejar la complejidad?
\end{center}
\begin{itemize}
 \item Cada función es un archivo
 \item No existe la modularidad
 \item Orientación a objetos inútil
 \item No existen los namespaces
 \item wrapper $eq$ mex
\end{itemize}
\end{frame}


\begin{frame}
\frametitle{Más sobre los wrappers}
\begin{itemize}
 \item Un wrapper es un adaptador entre los tipos de C, Fortran y los del lenguaje interpretado.
 \item En un lenguaje interpretado los tipos siempre son structs.
 \item En python los tipos tienen $\sim$ 30 elementos.
 \item El intérprete tiene su propio modelo de stack.
\end{itemize}
\end{frame}

\begin{frame}[containsverbatim]
 \frametitle{Arrays}
El objetivo es encajar
\begin{verbatim}
 double array[ ];
\end{verbatim}
en:
\begin{verbatim}
typedef struct PyArrayObject { 
    PyObject HEAD
    char * data;
    int nd;
    intp * dimensions;
    intp * strides;
    PyObject * base;
    PyArray Descr * descr;
    int flags;
    PyObject * weakreflist;
} PyArrayObject ; };
\end{verbatim}
\end{frame}

\begin{frame}
 \frametitle{Posibles problemas}
\begin{itemize}
 \item Hacerlo manualmente requiere C medio
 \item Conocimiento del intérprete
 \item Problemas de alineación (strides)
 \item ¿Como Fortran o como C?
\end{itemize}
\begin{center}
 ¿No era tan fácil?
\end{center}
\end{frame}

\begin{frame}
 \frametitle{Es fácil gracias a...}
\begin{itemize}
 \item ctypes
 \item f2py
\end{itemize}

\end{frame}

\begin{frame}
 \frametitle{ctypes}
Permite enlazar en tiempo de ejecución una librería al intérprete
\end{frame}

\begin{frame}[containsverbatim]
 \frametitle{Un wrapper inútil.}
\begin{verbatim}
from ctypes import c_int, POINTER #1
import numpy as np
from numpy.ctypeslib import load_library,ndpointer #1

def dgesv(N,A,B):
    A = np.asfortranarray(A.astype(np.float64)) #2
    B = np.asfortranarray(B.astype(np.float64))

    cN=c_int(N)
    NRHS=c_int(1)
    LDA=c_int(N)
    IPIV=(c_int * N)()
    LDB=c_int(N)
    INFO=c_int(1)

    lapack=load_library('liblapack.so','/usr/lib/')#3
...
\end{verbatim}
\end{frame}

\begin{frame}[containsverbatim]
 \begin{verbatim}
     lapack.dgesv_.argtypes=[POINTER(c_int),
             POINTER(c_int),
             ndpointer(dtype=np.float64,
             ndim=2,
             flags='FORTRAN'),
             POINTER(c_int), POINTER(c_int),
             ndpointer(dtype=np.float64,
                       ndim=2,
                       flags='FORTRAN'),
             POINTER(c_int),POINTER(c_int)]#4

    lapack.dgesv_(cN,NRHS,A,LDA,IPIV,B,LDB,INFO)#5
    return B

print dgesv(2,np.array([[1,2],[1,4]]),np.array([[1,3]]))
\end{verbatim}
\end{frame}

\begin{frame}
 \begin{center}
  \begin{huge}
   ¡No hay que programar en C!
  \end{huge}
 \end{center}
\vspace{2cm}
\begin{itemize}
 \item FORTRAN (trailing underscore)
 \item Conversión de arrays
 \item Llamadas por referencia
 \item Toda subrutina puede ser una librería, sólo hay que compilarla de otra manera.
 \item Velocidad de ejecución $\sim$ Fortran
 \item Reciclaje
\end{itemize}
\end{frame}

\begin{frame}
 \frametitle{f2py}
Es una aplicación que es capaz de entender la mayoría del código en Fortran y lo convierte automáticamente en un módulo de Python
\end{frame}


\begin{frame}[containsverbatim]
 \frametitle{Aún más fácil con f2py}
\begin{verbatim}
!t.f90
subroutine withCallback(a, b, ipar, rpar, callback)
  real a,b, rpar(*)
  integer ipar(*)
  external callback
  print*, 'The parameters are', a,b, ipar(:3), rpar(:3)
  call callback(rpar, ipar)
end subroutine withCallback

subroutine theCallback(rpar, ipar)
  real rpar(*)
  integer ipar(*)
  print*, 'Here the callback is called', ipar(:3), rpar(:3)
end subroutine theCallback
\end{verbatim}
\end{frame}

\begin{frame}[containsverbatim]
 \frametitle{Y funciona...}
\begin{verbatim}
$ f2py -c -m callback t.f90 --fcompiler=gnu95

>>> from numpy import *
>>> import callback
>>> ipar=array([4,5,6])
>>> rpar=array([1.,2.,3.])
>>> callback.withcallback(8,9,ipar,rpar,
...    callback.thecallback)
 The parameters are   8.000000       9.000000    4      5 
         6   1.000000       2.000000       3.000000
 Here the callback is called           4       5     6  
1.000000       2.000000       3.000000
>>>
\end{verbatim}
\end{frame}

\begin{frame}
 \begin{center}
 \begin{huge}
  Ahora sí parece más fácil.
 \end{huge}
\end{center}
\end{frame}

\begin{frame}
 \frametitle{Python en paralelo}
\begin{itemize}
 \item Se lanza Python como proceso
 \item La comunicación entre los intérpretes puede hacerse mediante MPI
 \begin{itemize}
  \item No hay wrappers para blacs pero pueden hacerse
 \end{itemize}
\end{itemize}
Por ejemplo:
\end{frame}

\begin{frame}
 \frametitle{GIL}
\begin{itemize}
 \item Cpython no es thread safe
\end{itemize}

\end{frame}


\end{document}
